\documentclass[10pt,a4paper]{amsart}
\usepackage{amsmath,amssymb,amsthm,graphicx}
% For commutative diagrams we use Xy-pic
\usepackage[all]{xy}

% We use 2cell for 2-commutative diagrams.
\xyoption{2cell}
% Theorem environments.
%
\theoremstyle{plain}
\newtheorem{theorem}[subsection]{Theorem}
\newtheorem{proposition}[subsection]{Proposition}
\newtheorem{lemma}[subsection]{Lemma}

\theoremstyle{definition}
\newtheorem{definition}[subsection]{Definition}
\newtheorem{example}[subsection]{Example}
\newtheorem{exercise}[subsection]{Exercise}
\newtheorem{situation}[subsection]{Situation}

\theoremstyle{remark}
\newtheorem{remark}[subsection]{Remark}
\newtheorem{remarks}[subsection]{Remarks}

\numberwithin{equation}{subsection}

% Macros
%
\def\lim{\mathop{\mathrm{lim}}\nolimits}
\def\colim{\mathop{\mathrm{colim}}\nolimits}
\def\Spec{\mathop{\mathrm{Spec}}}
\def\Hom{\mathop{\mathrm{Hom}}\nolimits}
\def\Ext{\mathop{\mathrm{Ext}}\nolimits}
\def\SheafHom{\mathop{\mathcal{H}\!\mathit{om}}\nolimits}
\def\SheafExt{\mathop{\mathcal{E}\!\mathit{xt}}\nolimits}
\def\Sch{\mathit{Sch}}
\def\Mor{\mathop{\mathrm{Mor}}\nolimits}
\def\Ob{\mathop{\mathrm{Ob}}\nolimits}
\def\Sh{\mathop{\mathit{Sh}}\nolimits}
\def\NL{\mathop{N\!L}\nolimits}
\def\CH{\mathop{\mathrm{CH}}\nolimits}
\def\proetale{{pro\text{-}\acute{e}tale}}
\def\etale{{\acute{e}tale}}
\def\QCoh{\mathit{QCoh}}
\def\Ker{\mathop{\mathrm{Ker}}}
\def\Im{\mathop{\mathrm{Im}}}
\def\Coker{\mathop{\mathrm{Coker}}}
\def\Coim{\mathop{\mathrm{Coim}}}

% Boxtimes
%
\DeclareMathSymbol{\boxtimes}{\mathbin}{AMSa}{"02}

%
% Macros for moduli stacks/spaces
%
\def\QCohstack{\mathcal{QC}\!\mathit{oh}}
\def\Cohstack{\mathcal{C}\!\mathit{oh}}
\def\Spacesstack{\mathcal{S}\!\mathit{paces}}
\def\Quotfunctor{\mathrm{Quot}}
\def\Hilbfunctor{\mathrm{Hilb}}
\def\Curvesstack{\mathcal{C}\!\mathit{urves}}
\def\Polarizedstack{\mathcal{P}\!\mathit{olarized}}
\def\Complexesstack{\mathcal{C}\!\mathit{omplexes}}
% \Pic is the operator that assigns to X its picard group, usage \Pic(X)
% \Picardstack_{X/B} denotes the Picard stack of X over B
% \Picardfunctor_{X/B} denotes the Picard functor of X over B
\def\Pic{\mathop{\mathrm{Pic}}\nolimits}
\def\Picardstack{\mathcal{P}\!\mathit{ic}}
\def\Picardfunctor{\mathrm{Pic}}
\def\Deformationcategory{\mathcal{D}\!\mathit{ef}}
\usepackage{cleveref}
\crefname{originalreference}{original reference}{original references}
\Crefname{originalreference}{Original reference}{Original references}
\makeatletter
\newlabel{lemma-blow-up-square-cohomological-descent}{{0EW5}{source}{}{sourcepointer-1}{}}
\newlabel{lemma-blow-up-square-cohomological-descent@cref}{{[originalreference][0][]0EW5}{[1][0][]source}{}{}{}}
\makeatother
\crefname{originalsection}{original section}{original sections}
\Crefname{originalsection}{Original section}{Original sections}
\makeatletter
\newlabel{sites-cohomology-section-leray}{{072X}{source}{}{sourcepointer-2}{}}
\newlabel{sites-cohomology-section-leray@cref}{{[originalsection][0][]072X}{[1][0][]source}{}{}{}}
\makeatother
\title{The blow-up square yields a long exact cohomology sequence for torsion coefficients}
\author{The Stacks Project Authors}
\date{stacks-project:etale-cohomology.tex:local-snapshot-20261011; GNU FDL 1.2 or later}
\begin{document}
\maketitle
\noindent\textit{Editorial note.} All delivered mathematical text is the author's own, copied from the pinned Stacks Project chapter file etale-cohomology.tex; no statement or proof step was rewritten, reordered or removed.\par\medskip
\noindent\textit{Editorial note.} The counted proof is the author's complete argument: the proof is written out in full: the descent triangle and the identification of the cohomology of the three terms through the section formula for a pushforward.\par\medskip
\noindent\textit{Editorial note.} The label lemma-blow-up-square-etale-cohomology names the counted statement, and the proof environment immediately following it in the source is delivered verbatim as its proof.\par\medskip
\noindent\textit{Editorial note.} The counted statement is carried by the author's own lemma environment in the chapter source.\par\medskip
\noindent\textit{Editorial note.} The statement's blow-up square is the author's and is delivered verbatim; the two cited inputs are kept as external pointers with their canonical Stacks tags.\par\medskip
\noindent\textit{Allowed standard prerequisites (editorial).}\begin{itemize}
\item The cohomological descent triangle of the blow-up square. Same chapter, the lemma producing the distinguished triangle K -> Ri\_*(K|\_Z) + Rb\_*(K|\_\{X'\}) -> Rc\_*(K|\_E) -> K[1]. Supplies the triangle whose long exact cohomology sequence is the assertion.
\item Derived global sections of a pushforward. Cohomology on Sites, the Leray section describing R Gamma(X, Rb\_* of a complex) as R Gamma(X' , -). Supplies the identification of the cohomology groups of the three terms, and similarly for the other pushforwards.
\end{itemize}\medskip
\ifdefined\hypertarget\hypertarget{sourcepointer-1}{}\fi
\noindent\textit{Source reference.} Original reference 0EW5: sources/stacks-etale-cohomology/pinned-original-source.tex, source lines 23315--23315. Prerequisite pointer invoked at proof line(s) 23411; it names a standard fact of the same chapter that the counted argument uses, and every step of the counted proof built on it is retained verbatim in the delivered body.\par\medskip
\ifdefined\hypertarget\hypertarget{sourcepointer-2}{}\fi
\noindent\textit{Source reference.} Original section 072X: sources/stacks-chow/tags.txt, source lines 9041--9041. Prerequisite pointer invoked at proof line(s) 23416; it names a standard fact of a sibling chapter that the counted argument uses, and every step of the counted proof built on it is retained verbatim in the delivered body.\par\medskip
\begin{lemma}
\label{lemma-blow-up-square-etale-cohomology}
Let $X$ be a scheme and let $K \in D^+(X_\etale)$ have
torsion cohomology sheaves. Let $Z \subset X$ be a closed subscheme
cut out by a quasi-coherent ideal of finite type. Consider the
corresponding blow up square
$$
\xymatrix{
E \ar[d] \ar[r] & X' \ar[d]^b \\
Z \ar[r] & X
}
$$
Then there is a canonical long exact sequence
$$
H^p_\etale(X, K) \to
H^p_\etale(X', K|_{X'}) \oplus
H^p_\etale(Z, K|_Z) \to
H^p_\etale(E, K|_E) \to
H^{p + 1}_\etale(X, K)
$$
\end{lemma}

\begin{proof}[First proof]
This follows immediately from
Lemma \ref{lemma-blow-up-square-cohomological-descent}
and the fact that
$$
R\Gamma(X, Rb_*(K|_{X'})) = R\Gamma(X', K|_{X'})
$$
(see Cohomology on Sites, Section \ref{sites-cohomology-section-leray})
and similarly for the others.
\end{proof}

\end{document}
